\documentclass{standalone}
\usepackage{circuitikz}
\usetikzlibrary{calc}
\definecolor{circuitnotemark}{HTML}{FE00FE}
\begin{document}
\begin{circuitikz}[american, line width=0.8pt]
\ctikzset{bipoles/length=1.2cm}
\foreach \x in {0,2,4,6,8,10} {\foreach \y in {0,-2,-4,-6} {\fill[gray, opacity=0.35] (\x,\y) circle (1.1pt);}}
\node[gray, font=\scriptsize] at (0,0.84) {1};
\node[gray, font=\scriptsize] at (2,0.84) {2};
\node[gray, font=\scriptsize] at (4,0.84) {3};
\node[gray, font=\scriptsize] at (6,0.84) {4};
\node[gray, font=\scriptsize] at (8,0.84) {5};
\node[gray, font=\scriptsize] at (10,0.84) {6};
\node[gray, font=\scriptsize, anchor=east] at (-0.84,0) {a};
\node[gray, font=\scriptsize, anchor=east] at (-0.84,-2) {b};
\node[gray, font=\scriptsize, anchor=east] at (-0.84,-4) {c};
\node[gray, font=\scriptsize, anchor=east] at (-0.84,-6) {d};
\coordinate (x1y1) at (0,0);
\coordinate (x2y1) at (2,0);
\coordinate (x5y3) at (8,-4);
\coordinate (x6y3) at (10,-4);
\coordinate (x1y4) at (0,-6);
\coordinate (x2y4) at (2,-6);
\coordinate (x5y2) at (8,-2);
\coordinate (x6y2) at (10,-2);
\coordinate (x5y1) at (8,0);
\coordinate (x2y2) at (2,-2);
\draw (x1y1) to[R, l_=$R_{1}$] (x2y1); % line 3
\draw (x5y3) to[R, l_=$R_{2}$] (x6y3); % line 4
\draw (x1y4) to[R, l_=$R_{3}$] (x2y4); % line 5
\draw (x5y2) to[R, l_=$R_{4}$] (x6y2); % line 6
\draw (x2y1) -| (x5y3); % line 8
\draw (x2y4) |- (x5y2); % line 9
\node[circ] at (x5y2) {};
\path (12,-7.408) rectangle (17.774,0.593); % line 11
\node[anchor=west, inner sep=0, circuitnotemark, font=\large] at (12,0) {X}; % line 11
\node[anchor=west, inner sep=0, circuitnotemark, font=\large] at (12,-0.487) {X}; % line 11
\node[anchor=west, inner sep=0, circuitnotemark, font=\large] at (12,-0.974) {X}; % line 11
\node[anchor=west, inner sep=0, circuitnotemark, font=\large] at (12,-1.46) {X}; % line 11
\node[anchor=west, inner sep=0, circuitnotemark, font=\large] at (12,-1.947) {X}; % line 11
\node[anchor=west, inner sep=0, circuitnotemark, font=\large] at (12,-2.434) {X}; % line 11
\node[anchor=west, inner sep=0, circuitnotemark, font=\large] at (12,-2.921) {X}; % line 11
\node[anchor=west, inner sep=0, circuitnotemark, font=\large] at (12,-3.408) {X}; % line 11
\node[anchor=west, inner sep=0, circuitnotemark, font=\large] at (12,-3.895) {X}; % line 11
\node[anchor=west, inner sep=0, circuitnotemark, font=\large] at (12,-4.381) {X}; % line 11
\node[anchor=west, inner sep=0, circuitnotemark, font=\large] at (12,-4.868) {X}; % line 11
\node[anchor=west, inner sep=0, circuitnotemark, font=\large] at (12,-5.355) {X}; % line 11
\node[anchor=west, inner sep=0, circuitnotemark, font=\large] at (12,-5.842) {X}; % line 11
\node[anchor=west, inner sep=0, circuitnotemark, font=\large] at (12,-6.329) {X}; % line 11
\node[anchor=west, inner sep=0, circuitnotemark, font=\large] at (12,-6.816) {X}; % line 11
\node[anchor=south west, inner sep=2pt, yshift=4pt, circuitnotemark, font=\LARGE\bfseries] (circuittitle) at (current bounding box.north west) {X};
\path (circuittitle.south west) rectangle ++(5.578,0.853);
\node[anchor=north east, inner sep=2pt, circuitnotemark, font=\large] (circuitstamp) at ([yshift=-0.593cm]current bounding box.south east) {X};
\path (circuitstamp.north east) rectangle ++(-5.08,-0.593);
\end{circuitikz}
\end{document}
